\subsection{Сложное движение твердого тела}
%%%%%%%%%%%%%%%%%%%%%%%%%%%%%%%%%%%%%%%%%%%%%%%%%%%%
%%%%%%%%%%%%%%%%%%%%%%%%%%%%%%%%%%%%%%%%%%%%%%%%%%%% 
\z[{\begin{wrapfigure}[4]{r}{0.2\textwidth}%
		\vspace{-10pt}%
		\begin{flushright}%
 			\includegraphics[width=0.2\textwidth]{8.4.1}%
		\end{flushright}%
	\end{wrapfigure}%
	}]%
	{На горизонтальной плоскости лежит катушка ниток. Катушку тянут за нитку c силй $F$. При каких углах $\alpha$ между силой и горизонталью катушка станет ускоряться в сторону натянутой нити? Чему равно это ускорение?}
%
{$\cos\alpha > \frac{r}{R}$}
%%%%%%%%%%%%%%%%%%%%%%%%%%%%%%%%%%%%%%%%%%%%%%%%%%%%	
\z[{\begin{wrapfigure}[4]{r}{0.2\textwidth}%
		\vspace{-10pt}%
		\begin{flushright}%
 			\includegraphics[width=0.2\textwidth]{8.4.2}%
		\end{flushright}%
	\end{wrapfigure}%
	}]%
	{Оси тонкостенного и сплошного цилиндров соединены невесомой штангой. Цилиндры скатываются без проскальзывания по наклонной плоскости с углом $\alpha$. Радиусы цилиндров одинаковы, масса каждого цилиндра $m$. Определите силу натяжения $T$ штанги.}
%
{$T = \frac{1}{7}mg\sin\alpha$}
%%%%%%%%%%%%%%%%%%%%%%%%%%%%%%%%%%%%%%%%%%%%%%%%%%%% 
\z[{\begin{wrapfigure}[4]{r}{0.2\textwidth}%
		\vspace{-10pt}%
		\begin{flushright}%
 			\includegraphics[width=0.2\textwidth]{8.4.3}%
		\end{flushright}%
	\end{wrapfigure}%
	}]%
	{На тонкостенный цилиндр намотана нить, конец которой закреплен на стойке так, что при соскальзывании цилиндра с наклонной плоскости нить остается параллельной наклонной плоскости. Какую скорость $v$ приобрел цилиндр, если его ось прошла расстояние $l$? Угол наклона плоскости $\alpha$, коэффициент трения между плоскостью и цилиндром $\mu$.}
%
{$v = \sqrt{gl(\sin\alpha - 2\mu\cos\alpha)}$}
%%%%%%%%%%%%%%%%%%%%%%%%%%%%%%%%%%%%%%%%%%%%%%%%%%%%
\z[{\begin{wrapfigure}[5]{r}{0.2\textwidth}%
		\vspace{-15pt}%
		\begin{flushright}%
 			\includegraphics[width=0.2\textwidth]{8.4.4}%
		\end{flushright}%
	\end{wrapfigure}%
	}]%
	{Тонкостенный цилиндр катится по горизонтальной поверхности без проскальзывания со скоростью $v_0 = 6\,\text{м/с}$. Коэффициент трения между цилиндром и поверхностью равен $\mu = 0.2$. Цилиндр сталкивается с вертикальной гладкой стенкой и упруго отражается от неё. Найдите: 
	\lpu скорости верхней и нижней точек цилиндра непосредственно после отражения;\lpuend
	\lpu скорость центра цилиндра через $t_1 = 2\,\text{с}$ после столкновения со стенкой и путь $s_1$, пройденный им за это время; \lpuend
	\lpu какой путь $s_2$ пройдет центр цилиндра к моменту времени $t_2 = 4\,\text{с}$. \lpuend
}
%
{\puan $v_B = 0$, $v_\text{н} = 2v_0$;~\puan $2\,\text{м/с}$ и $8\,\text{м}$;~\puan $9\,\text{м}$}
%%%%%%%%%%%%%%%%%%%%%%%%%%%%%%%%%%%%%%%%%%%%%%%%%%%%	
\z[{\begin{wrapfigure}[7]{r}{0.25\textwidth}%
		\vspace{-20pt}%
		\begin{flushright}%
 			\includegraphics[width=0.25\textwidth]{8.4.5}%
		\end{flushright}%
	\end{wrapfigure}%
	}]%
	{«Хула-хупп» это обруч, с которым можно проделывать фокусы. Например, закручивают его в вертикальной плоскости и толкают от себя по горизонтали, после чего вращающийся обруч, проскальзывая по полу, отъезжает от них на некоторое расстояние и возвращается обратно. Сформулируем задачу так: обруч радиусом $R$ в момент толчка закручен вокруг горизонтальной оси до угловой скорости $\omega_0$, и ему придали скорость $v_0$, направленную вдоль пола перпендикулярно оси вращения. Коэффициент трения обруча об пол равен $\mu$.
	\lpu Через какое время $t$ обруч начнет двигаться назад?\lpuend
	\lpu На какое максимальное расстояние $s_{\text{max}}$ в направлении скорости $v_0$ обруч удалится от начальной точки?\lpuend
	\lpu Какова будет его угловая скорость $\omega$ в момент остановки? \lpuend
	\lpu  С какой скоростью $v_{\text{кон}}$ он будет катиться после окончания проскальзывания по полу? \lpuend
}
%
{\puan $s_\text{макс} = \frac{v_0^2}{2\mu g}$;~\puan $\omega_1 = \omega_0 - \frac{v_0}{R}$;~\puan $v = \frac{\omega_0 R - v_0}{2}$;~\puan $v_0 = \frac{1}{3}\omega_0 R$ }